\appendix

\section{HRP Recovery Proof}
\label{app:proof}

We provide the full proof of Proposition~\ref{prop:recovery}.

\begin{proof}
Let the DHRP layer have a fixed binary tree of depth $L$ with internal nodes $\mathcal{K}$ and leaf nodes $\mathcal{L}$. Each internal node $k \in \mathcal{K}$ has a gating network $g_k: \mathbb{R}^d \to \mathbb{R}^2$ and routing probability $\pi_k(\tau) = \mathrm{softmax}(g_k(h_t)/\tau)$.

\textbf{Step 1: Pointwise convergence of gates.}
Fix any internal node $k$. Since $g_k$ has a unique maximizing direction, $g_k(h_t)[0] \neq g_k(h_t)[1]$. Without loss of generality, suppose $g_k(h_t)[0] > g_k(h_t)[1]$. Then:
\[
\pi_k(\tau)[0] = \frac{e^{g_k(h_t)[0]/\tau}}{e^{g_k(h_t)[0]/\tau} + e^{g_k(h_t)[1]/\tau}} = \frac{1}{1 + e^{(g_k(h_t)[1] - g_k(h_t)[0])/\tau}}.
\]
Since $g_k(h_t)[1] - g_k(h_t)[0] < 0$, as $\tau \to 0^+$, the exponent $\to -\infty$, so $\pi_k(\tau)[0] \to 1$ and $\pi_k(\tau)[1] \to 0$. This recovers the hard assignment $\pi_k^{\text{hard}} = (1, 0)$.

\textbf{Step 2: Convergence of leaf budgets.}
For each leaf $\ell \in \mathcal{L}$, the budget is $p_\ell(\tau) = \prod_{(k,d) \in P_\ell} \pi_k(\tau)[d]$, where $P_\ell$ is the root-to-leaf path. Since the tree has finite depth $L$, this is a finite product of $L$ terms. By Step~1, each factor converges pointwise to either 0 or 1. Therefore:
\[
\lim_{\tau \to 0^+} p_\ell(\tau) = \prod_{(k,d) \in P_\ell} \pi_k^{\text{hard}}[d] \in \{0, 1\}.
\]
Exactly one leaf in each subtree receives budget 1 (the leaf reached by following the argmax direction at each node), reproducing the hard recursive bisection of classical HRP.

\textbf{Step 3: Recovery of weights.}
With $\alpha = 1$, the blending equation reduces to $w_i = b_i$, where $b_i = p_{\ell(i)} / \sum_j p_{\ell(j)}$. In the $\tau \to 0^+$ limit, $b_i$ equals the HRP weight when the tree structure and inverse-variance allocation rule match, since each asset is assigned to exactly one leaf with the same risk budget split as HRP's recursive bisection.

The interchange of limit and finite product is justified because the product has a fixed finite number of factors (equal to tree depth $L$), and each factor converges individually. No dominated convergence argument is needed for the product; pointwise convergence of each factor suffices for a finite product.
\end{proof}

\section{Hyperparameter Configuration}
\label{app:hyperparams}

Table~\ref{tab:hyperparams} summarizes the hyperparameters used across all experiments.

\begin{table}[h]
\centering
\caption{Hyperparameter settings for DHRP and the LLM-DHRP$^{\dagger}$ negative ablation across all three universes. The LLM-DHRP$^{\dagger}$ rows are reported for full reproducibility of the negative result analysed in Section~\ref{sec:llmablation}.}
\label{tab:hyperparams}
\begin{tabular}{lrrr}
\toprule
Hyperparameter & DM & EM & CM \\
\midrule
Tree depth $L$ & 3 & 3 & 4 \\
Hidden dimension & 64 & 64 & 64 \\
Learning rate & $4.5 \times 10^{-4}$ & $4.5 \times 10^{-4}$ & $4.5 \times 10^{-4}$ \\
Weight decay & $3 \times 10^{-4}$ & $3 \times 10^{-4}$ & $3 \times 10^{-4}$ \\
Batch size & 32 & 32 & 32 \\
Training epochs & 40 & 40 & 40 \\
CRRA $\gamma$ & 2.5 & 1.2 & 2.5 \\
$\lambda_{\text{HRP}}$ schedule & $0.3 \to 0.1$ & $0.6 \to 0.3$ & $0.3 \to 0.1$ \\
Routing temperature $\tau$ & 1.0 & 1.0 & 1.0 \\
Training window (days) & 252 & 252 & 252 \\
Holding period (days) & 21 & 21 & 21 \\
Purge window (days) & 5 & 5 & 5 \\
Gradient clip norm & 1.0 & 1.0 & 1.0 \\
Optimizer & AdamW & AdamW & AdamW \\
\midrule
\multicolumn{4}{l}{\textit{LLM-DHRP additional parameters}} \\
FinBERT embedding dim & 768 & 768 & 768 \\
FinBERT projection dim & 64 & 64 & 64 \\
Qwen3 sentiment features & 16 & 16 & 16 \\
Cross-attention heads & 4 & 4 & 4 \\
Gate initialization & 0.0 & 0.0 & 0.0 \\
\bottomrule
\end{tabular}
\end{table}

\section{Reproducibility Statement}
\label{app:reproducibility}

All code, model architectures, training procedures, and evaluation pipelines are available in the accompanying repository.\footnote{Repository URL will be included in the camera-ready version.} The repository includes:
\begin{itemize}
\item Complete source code for DHRP and LLM-DHRP layers (\texttt{src/models/})
\item All baseline implementations including deep baselines (\texttt{src/models/deep\_baselines.py})
\item Data loading and feature engineering pipelines (\texttt{src/data/})
\item Evaluation framework including backtesting, statistical tests, and visualization (\texttt{src/evaluation/})
\item Pre-trained model checkpoints (\texttt{results/models/})
\item Jupyter notebook reproducing all experiments (\texttt{notebooks/llm\_dhrp\_experiments.ipynb})
\end{itemize}

Price data are sourced from Yahoo Finance (freely available). Text data are sourced from GDELT and Google News RSS (publicly accessible). Fama-French factors are from Kenneth French's data library. All experiments were run on a single NVIDIA A100 GPU (40~GB) using PyTorch 2.x. Total training time for all models across all three universes is approximately 30 minutes. FinBERT inference requires approximately 2 minutes on a T4 GPU; Qwen3-8B inference requires approximately 5 minutes on an A100.

\section{Broader Impact Statement}
\label{app:impact}

This work presents a methodological contribution to portfolio optimization. The DHRP layer is a general-purpose tool for embedding hierarchical risk diversification within differentiable pipelines. We do not foresee direct negative societal consequences from this work; however, as with all financial modeling tools, practitioners should exercise caution when deploying learned allocation strategies in live markets. Model outputs should be treated as decision-support signals rather than autonomous trading directives. The ten-year backtest period, while covering multiple regimes, represents a single historical realization. Overfitting to past market dynamics is a pervasive risk in quantitative finance that no backtest methodology can fully eliminate.

\section{Correlation Structure}

\begin{figure}[h]
\centering
\includegraphics[width=0.85\linewidth]{correlations.png}
\caption{Asset correlation matrices. Developed markets exhibit diverse correlations enabling hierarchical diversification; emerging markets are uniformly positive ($\bar{\rho} \approx 0.55$).}
\label{fig:correlations_app}
\end{figure}

\section{Return Distributions}

\begin{figure}[h]
\centering
\includegraphics[width=0.95\linewidth]{distributions.png}
\caption{Return distributions by method and market. DHRP maintains favorable tail properties while shifting the distribution rightward.}
\label{fig:distributions_app}
\end{figure}
